\chapter{Sub-monolayer Tracking of Molecular States at Buried Interfaces}\label{chap:znpc}

Prepared for submission to an academic journal.

Harlan Heilman, Em Biegler, Sushil Adhikari, Victor Murcia, Tanner Melody, Fred Woodall, Lars Thomsen, Obaid Alqahtani, and Brian A. Collins, ``Sub-monolayer tracking of molecular states at buried interfaces,'' arXiv:2610.31924 (2026).\autocite{heilman2026sub}

This chapter is adapted from the manuscript above, which is under peer review and posted as a preprint on arXiv. The first question of \cref{chap:introduction} is examined here by applying the film tensor of \cref{chap:molecular-films} and the transfer matrix forward model and fitting workflow of \cref{chap:methods} to polarized resonant reflectivity of ZnPc films on silicon. Figure and table numbers follow the dissertation, with the supporting information in Supplemental Material.

\section{Attributions}\label{sec:attributions}

HH and BAC conceived the study. BAC oversaw and directed all aspects of the study. HH performed RSoXS measurements, DFT calculations, developed and performed the structural and optical modeling, and wrote the manuscript. FW deposited samples measured with RSoXS and performed and measured AFM topology. EB deposited a separate sample, collected, and analyzed non-resonant XRR. SA performed and analyzed VASE. TM collected and analyzed GIWAXS and measured AFM topography. VM assisted with DFT calculations and clustering. LT collected AR-NEXAFS and OA processed and analyzed the AR-NEXAFS.

\section{Introduction}\label{sec:introduction}

The molecular orientation\autocite{tumbleston2014influence,yokoyama2011molecular}, composition\autocite{collins2013absolute} and oxidation state\autocite{cabana2018mechanisms,wang2024insights} at the inaccessible interfaces of buried nanoscale assemblies are decisive quantities that dictate material function across physics, chemistry and materials science critical to advanced technologies. For example, in extreme-ultraviolet photoresists the chemical gradient at the edge of the latent image limits the minimum achievable process size for microprocessors.\autocite{debisschop2017stochastic,andrle2026nanostructure} In organic devices such as solar cells, LEDs, transistors, and spin valves the orientation and conformation of the molecular monolayer at the interface governs emergent electronic states that control charge generation, recombination, transport, and spin hybridization respectively.\autocite{tumbleston2014influence,yokoyama2011molecular,pei2020semiconductor,jiang2025toughened,droghetti2016dynamic} In bio- and chemical sensing, molecular states at buried interfaces have resulted in novel ionic transport dynamics.\autocite{khan2025local,trigg2018self,dong2010super} Controlling molecular states directly at these interfaces is, therefore, central to advancing each of these technologies.

Unfortunately, the few techniques that resolve buried interfacial structure, such as atom probe tomography are laborious, highly destructive, and can't resolve details of bonding states and their orientations.\autocite{cojocarumiredin2024atom} Lacking a direct probe, novel interfacial molecular states (e.g. orientation and packing) are thought to be achieved via substrate patterning effects,\autocite{kato2020morphological,ngongangndjawa2015impact,schunemann2011zinc} similar to atomic epitaxial growth. Studies have attempted to access the buried molecular state by measuring deposition at coverage checkpoints\autocite{derrien2022interfacial,yamamoto2022variation,deoteyza2006structural} with in situ measurements resolving a density stratification that contradict the patterning hypothesis.\autocite{derrien2022interfacial} Yet other work suggests these coverage snapshots only describe the exposed molecular layers in a transient state that change once buried,\autocite{yamamoto2022variation,deoteyza2006structural} leaving the mechanisms governing formation and control of interfacial states unresolved.

RSoXS, tuned to the energies of electronic transitions into unoccupied molecular orbitals can non-destructively reveal the chemical, electronic, and physical states within these buried structures\autocite{grabner2025resonant,andrle2026nanostructure,andrle2025using}. The scattered intensity is determined by two coupled effects, the geometric structure of the assembly and the optical interaction with the material. For reflectivity analysis, models of these effects are fit to scattering data to reveal depth-dependent information.\autocite{grabner2025resonant,ferron2022characterization,wang2007resonant,yan2013accurate,mezger2011molecular,ferron2023composition} Unfortunately, the physio-chemical states must be inferred from of optical parameters as a function of energy and polarization which increase in number for each energy probed, causing the analysis to become unstable. Therefore, analyses retreat to a single diagnostic resonance which cannot distinguish chemical species nor orientation unless simplifying assumptions are made.\autocite{ferron2022characterization,li2026pattern,nguyen2024resonant} Worse, a buried interfacial layer is typically required to accurately fit the data which is assumed to be comprised of contaminants or oxidation,\autocite{ferron2022characterization,wang2007resonant,mezger2011molecular,ferron2023composition} and there is no way to determine if these layers are real or merely a limitation of the soft X-rays' diffraction limit which sets the smallest resolvable feature by the scatterer's wavelength, in this case, several nanometers.

Recently, near edge X-ray absorption fine structure (NEXAFS) spectroscopy measurements have been used to construct an Edisonian optical model, enabling the orientation state to be inferred though a secondary analysis of the optical parameters, although the above limitations remain.\autocite{ferron2022characterization,wang2007resonant,yan2013accurate,ferron2023composition,thelen2020molecular} Alternative approaches arrive at the optical model from density functional theory (DFT), yet are limited by fundamental computational inaccuracies.\autocite{pasquali2014structural,capelli2019quantitative,pasquali2021analysis} Recently, we have demonstrated that ML-refined DFT optical models can quantitatively reproduce angle-resolved (AR) NEXAFS measurements of polycrystalline thin films with different orientation and crystal phases while retaining individual bond identity and orientation.\autocite{murcia2026quantitative}

Here we construct and apply an ML-DFT optical model in RSoXS analysis to directly quantify molecular chemical identity, orientation and density as a function of depth in a molecular thin film. The approach ties spectral resonances together in one physical model of the molecule, to increase chemical and orientational sensitivity, and locks scattering features between energies and polarization to recover lost diffraction limited information. This results in sub-monolayer resolution of molecular chemical andorientational state. Analysis of the bottom interface reveals not contaminants or a substrate patterned crystal phase, but a smooth variation of molecular state with partial coverage of the molecule in a unique orientation and packing state directly at the heterointerface. This discovery establishes the bulk texture is set not by substrate patterning but by competing surface and intermolecular interactions during film growth. The new analysis, therefore, enables direct characterization of buried structures at the molecular level critical to the physics of crystal growth, kinetic deposition, and solvent interaction models that dictate these interfaces across material systems.

\begin{figure}
  \centering
  \includegraphics[width=0.45\linewidth,height=0.75\textheight,keepaspectratio]{fig4-01_ml-dft-tensor}
  \caption[An ML-DFT resonant optical model for ZnPc.]{An ML-DFT resonant optical model for ZnPc. ML-DFT (Traces) and angle resolved NEXAFS measurements (points), deviating by \(\epsilon = 2.7\%\) and \(\epsilon_{\max} = 5.3\%\) at resonance (1). Top, schematic of the molecular frame and the four inequivalent carbon excitation centers, colored as labeled. (a) \(\beta_{xx}\), \(\beta_{zz}\), and the isotropic average \(\beta\) (gray), numbered bars are fine structure resonances targeted for RSoXS measurement, with colored pips indicating the excitation centers dominating the transition. (b) \(\delta\) by Kramers-Kronig transformation,\autocite{kramers1927diffusion} with shaded \(\Delta\delta\) bounded by the zero crossings \(E^{0}_{xx}\) and \(E^{0}_{zz}\). Vertical purple dash is a diagnostic RSoXS energy.}
  \label{fig:4.1}
\end{figure}


\section{Constructing an ML-Refined DFT Optical Tensor}\label{sec:constructing-an-ml-refined-dft-optical-tensor}

We focus on thin films of Zinc phthalocyanine (ZnPc), a mature molecular platform spanning foundational studies of single-molecule photophysics and quantum optics\autocite{yang2020sub,zhang2016visualizing,luo2019electrically,cao2021energy,doppagne2018electrofluorochromism}, commercial photovoltaics\autocite{klenkler2015photovoltaic,pfeiffer2013organic}, thin-film transistors and sensors\autocite{boileau2019metal,comeau2022surface}, and photomedicines.\autocite{chornovolenko2025phthalocyanines,rak2023advances,luo2021nanoscale} Physical vapor deposition yields polycrystalline films with the molecular orientation commonly set by patterning underlayers. CuI and graphene pattern face-on orientation (molecules lying flat) with brickstone packing while Si patterns the edge-on orientation with herringbone packing\autocite{kato2020morphological,ngongangndjawa2015impact,schunemann2011zinc}. However, in situ coverage studies have suggested that ZnPc and other metal-Pc films might bear more complex microstructural states that evolve from the substrate to the vacuum interface.\autocite{derrien2022interfacial,yang2015molecular,yamamoto2022variation,deoteyza2006structural,osso2004structure}

On silicon, ZnPc assumes a monoclinic herringbone crystal structure\autocite{schunemann2011zinc,derrien2022interfacial} with a strong edge-on texture confirmed by ellipsometry and GIXRD (Supplemental Material, \suppsec{1}). AR-NEXAFS places the film-averaged macrocycle normal (\(\widehat{z}\) in \cref{fig:4.1} schematic) at a polar tilt of \(\gamma = \ang{74}\) from the substrate surface normal (Supplemental Material, \suppsec{1}, i.e. edge-on orientation). The molecular index of refraction tensor \(n_{ij}(E) = 1 - \delta_{ij}(E) + i\beta_{ij}(E)\) follows by projecting the experimentally measured \(\beta_{\theta}(E)\) (from AR-NEXAFS) onto the molecular optical axes (\cref{fig:4.1}a, symbols)\autocite{stohr2010nexafs,collins2012polarized}, using the molecule-to-film relation given in Supplemental Material, \suppsec{2}. Unfortunately, the measured optical anisotropy is a convolution of the transition dipole moments (TDM) and the film-averaged molecular tilt. Through our ML clustering algorithm we remove this ambiguity to extract the orientation of the principle TDMs calculated from DFT simulated NEXAFS.\autocite{murcia2026quantitative} TDM clusters are then refined to quantitatively reproduce AR-NEXAFS (Supplemental Material, \suppsec{3}) compared in \cref{fig:4.1}. The average error between the refined optical tensor in reproducing the AR-NEXAFS data is 2.7\% with all experimental features reproduced in \cref{fig:4.1}. Thus, the ML-DFT single molecule tensor fully captures the optical interactions emergent in the film, while assigning chemical character and TDMs to spectral features.

The first two TDM clusters orient along \(\widehat{z}\) and carry the dominant oscillator strength, producing a strong pre-edge birefringence (\(\Delta\delta\), \cref{fig:4.1}b) that eventually flips at higher energies. We exploit this in reflectivity measurements via two modes. First, with high-density sampling in energy up to just below the first absorption resonance, we locate the extinction of the s-polarized critical angle \(\theta_{c}^{s} = \sqrt{2\delta_{xx}}\) to inform the film averaged orientation at low radiation dose. Maximum q-range reflectivity is then collected at two book-ending diagnostic energies, maximizing sensitivity to thin buried layers. This pre-edge data reports the orientation of a generic conjugated framework rather than chemical identity. Chemical specificity is provided though RSoXS measurements at the fine structure resonances which resolve all four inequivalentcarbon environments of the macrocycle (labeled in \cref{fig:4.1} schematic and circles in \cref{fig:4.1}a) as determined through the DFT calculation.

\begin{figure}
  \centering
  \includegraphics[width=\linewidth,height=0.75\textheight,keepaspectratio]{fig4-02_rsoxs-fits}
  \caption[Fitting energy and polarization-dependent RSoXS.]{Fitting energy and polarization-dependent RSoXS. (a) Reflectivity model description of the film using three ZnPc slabs on SiO\textsubscript{2}/Si with thickness \(t\) and roughness \(\sigma\) each, differing only by the optical parameterization of each slab with the Free model requiring four uniaxial components per energy per slab against the ML-DFT model with just two, molecular tilt \(\gamma\) and density \(\rho\) per slab. (b) Film out of plane tensor components built from \cref{fig:4.1} parameterized by tilt angle from face to edge on with the isotropic condition in black. Vertical lines denote the energies in the resonant series. (c-e) Reflectivity data (showing every fifth datapoint, Supplemental Material, \suppsec{5}), s-pol (filled circle) and p-pol (triangle, lighter), against the Free (solid) and ML-DFT (dashed) models for each energy domain: (b) diagnostic energies, (c) pre-edge, δ-crossing, (d) post-edge, (e) resonant. Ticks in (c) mark the Q-ranges of the largest divergences. Traced in (d) plots the critical angle for each s-and-p-pol across the pre-edge energies. Molecule frames inset in (e) highlight the carbon excitation centers within one outer ZnPc wing that dominate each TDM cluster.}
  \label{fig:4.2}
\end{figure}


\section{Direct, Quantitative Sensing of Molecular State with Depth}\label{sec:direct-quantitative-sensing-of-molecular-state-with-depth}

We analyze the reflectivity data across these energies via forward simulation where the film is modeled in a three-slab structure, each with a thickness \(t\) and roughness \(\sigma_{r}\), outlined in \cref{fig:4.2}a. Three slabs are required for a good fit and are presumed to capture unique chemical states and species at the surface, bulk and buried interface of the film. We then test two optical models to contrast the ML-DFT optical model with that of the traditional analysis. The traditional approach, designated the Free model, fits four uniaxial tensor components per energy per slab (\cref{fig:4.2}a). A secondary analysis must be developed to determine chemical nature, orientation, and density from these optical parameters. In contrast to this, the ML-DFT model uses the energy and polarization-dependent optical tensor from \cref{fig:4.1} (derived in Supplemental Material, \suppsec{4}, orientation dependence in \cref{fig:4.2}b) to directly parameterize density \(\rho\) and molecular tilt angle \(\gamma\) -- only two parameters total per slab and consistent across energy compared with 84 in the Free model. In this approach, different chemistries, such as oxidation or contamination in the surface, bulk or substrate layers will be manifest as fit variances at the molecule's key resonant energies (\cref{fig:4.2}e).

Fit quality between the two models differs by energy regime. Across the pre-edge and resonance series of \cref{fig:4.2}d and~e, the DFT model tracks the data with similar or better residuals than the free tensor model. The critical angles traced in \cref{fig:4.2}d, vertical black traces, highlight the underlying birefringence that is consistent with the ZnPc molecules oriented edge-on to the substrate. By contrast, the molecule's resonance series (\cref{fig:4.2}e) is sensitive to chemical and electronic orbital identity. Here the ML-DFT model fits better than the Free model (average deviation 1.7\(\sigma\) vs 2.2\(\sigma\), respectively). Since the ML-DFT model restricts only ZnPc to be present in the film, this result indicates a lack of measurable contamination or oxidation in the film to which interfacial layers are customarily assigned.\autocite{ferron2022characterization,mezger2011molecular,ferron2023composition,pasquali2014structural} Instead, these layers depict a unique molecular configurational state well described by the ML-DFT model.

\Cref{fig:4.2}c presents the largest divergences between the models, where the characteristics of the \qty{283.7}{eV} reflectivity profile are the most challenging to reproduce for the either model. At this energy, the ML-DFT model performs worse in the s-pol profile with a noticeable phase shift (ticks, \cref{fig:4.2}c) between the fringes in the data and model. Conversely, in the p-profile, the Free model fits worse and fails to capture the fringe structure around Brewster's angle (ticks, \cref{fig:4.2}c). Despite this, both models arrive at similar slab structures (\textbf{Tab.} \textbf{SI, SII}) that agree with total thicknesses from ellipsometry analysis (\suppfig{1}), yet the DFT model achieves this with more than \(10\times\) sharper parameter uncertainties than the Free model.

At this point, the Free model optical parameters must be analyzed further to interpret orientation and density. In contrast, these critical parameters are read directly from the ML-DFT model, quantitativelytracing the microstructure evolution through the film in three distinct densities and orientation states. The parameters of the central bulk layer agree well with tabulated density and orientation of herringbone ZnPc.\autocite{kato2020morphological,derrien2022interfacial} Separately, the surface and interfacial layers exhibit orientation and densities that do not agree with herringbone packing yet both carry the resonant anisotropy that only orientated ZnPc can produce -- in particular, the resonant amplitudes and even anisotropy flipping across the edge.

\begin{figure}
  \centering
  \includegraphics[width=0.45\linewidth,height=0.75\textheight,keepaspectratio]{fig4-03_model-comparison}
  \caption[Dominance of the ML-DFT model over traditional analyses.]{Dominance of the ML-DFT model over traditional analyses. (a) Birefringence \(\Delta\delta = \delta_{xx} - \delta_{zz}\) (circles, solid) and dichroism \(\Delta\beta = \beta_{xx} - \beta_{zz}\) (triangles, dashed) for the surface, bulk, and interfacial ZnPc layers. Orange points are Free tensor optical parameters fit independently per energy. Blue traces follow from the ML-DFT optical model with blue areas representing uncertainties. (b) Model selection criterion (MSC): reduced total parameter counts \(N\), \(\chi^{2}\), and the AIC and BIC criteria for the two models where lower values indicate the preferred model. Hashed regions in \(N\) represent constrained experimental offset parameters (Supplemental Material, \suppsec{5})}
  \label{fig:4.3}
\end{figure}


The optical properties of the two models, as shown in \cref{fig:4.3}a (expanded in \suppfig{10}), agree closely within the bulk, while those of the Free model surface and buried interface layers are both more extreme and highly uncertain, enough such that there are no significant disagreements between the models. The comparison of the models highlights several issues with the Free model. For instance, the unphysical dichroism at \qty{287}{eV} (black arrow pointing off scale in \cref{fig:4.3}a) highlights the fact that there are no constraints against unphysical results. In particular, there is no known way to assure the Free model optical parameters obey Kramers-Kronig consistency without excessive energy sampling.\autocite{ferron2023composition,stone2012kramers} The second notable issue is interpreting the Free model optical parameters for orientational or chemical state. In particular, the low energy trends in \cref{fig:4.3}a are discontinuous at the diagnostic dataset at \qty{283.7}{eV} (circled points in \cref{fig:4.3}a), resulting in contradicting conclusions of the molecular orientation state depending on which energy is analyzed.

Not only is the ML-DFT model superior, being physically constrained and interpretable, but all statistical metrics demonstrate its advantages (\cref{fig:4.3}b) due especially to the reduced total number of parameters in the model (\emph{N}). The reduced chi square parameter \(\chi_{\nu}^{2}\) suggests that the DFT model reaches the best agreement with the experiment across the full set of energies. The Bayesian and Akaike information criterion (BIC \& AIC) measure the weight of evidence for each model and the expected performance for new measurements, respectively. From these results, we conclude that the DFT model achieves a drastic statistical improvement in fit quality with \(17\times\) fewer parameters and \(10\times\) smaller parameter uncertainties, while presenting directly interpretable results that are useful to interrogate the molecular state.

\begin{figure}
  \centering
  \includegraphics[width=\linewidth,height=0.75\textheight,keepaspectratio]{fig4-04_depth-profile}
  \caption[Revealing sub-monolayer orientation and density gradients at a buried interface.]{Revealing sub-monolayer orientation and density gradients at a buried interface. (a) s- and p-pol reflectivity fits to \qty{283.7}{eV} and resonant energies for the slab model (dashed), against the orientation gradient (red). For clarity, the q-range is stretched by \(2\times\) for the resonant profiles. Inset reports model fit quality as \(2\sigma\). (b-d) Real space analysis of molecular state versus depth. (b) orientation and (d) density depth profile of the gradient model with characteristic interfacial widths \(\tau_{\mathrm{Si}}\) and \(\tau_{\mathrm{vac}}\) and RMS surface roughness, shaded bands representing 95\% confidence (2σ). (c) Visual reconstruction of the molecular state with depth consistent with orientation, density, and diffraction measurements. Gray molecules represent a layer behind the colored molecular layer in the front.}
  \label{fig:4.4}
\end{figure}


\section{Sub-monolayer Interfacial Physics}\label{sec:sub-monolayer-interfacial-physics}

Without needing to account for unique chemical species at the boundaries of the ZnPc bulk layer, the independent slab model is no longer physically motivated. Instead, we develop a single slab gradient model based on the physics of a continuous second order transition from a pinned molecular state due to surface energies that evolve toward an energy-minimized state in the bulk (derivation in \cref{sec:orientation-gradient-theoretical-foundation}). This results in a description physically matched to an evolving growth front during deposition with each deposited layer taking its orientation from those beneath it. This gradient model resolves orientation as a smooth function of depth, parameterized by the bulk state and two exponentials that transition over characteristic lengths \(\tau\) toward the pinned state at each interface. Besides incorporating the physics at play during deposition, this structural model also has fewer parameters than the slab model.

The fit using the graded interface model (with ML-DFT modeling optics) is compared to the slab model in \cref{fig:4.4}a for the diagnostic profile at \qty{283.7}{eV}, together with the resonant series but neglecting the rest of the data. Notably, the erroneous phase shift of the fringes in the slab model is removed in the gradient model and the Brewster angle is reproduced significantly better than the slab model. These fit improvements along with the reduced number of parameters enable the best goodness of fit, indicating this model recovers the true orientation and density gradient within the sample more accurately than what the traditional slab picture can achieve.

\begin{table}
  \centering
  \caption{Best fit parameters from the orientation and density gradient model. Total film thickness is \(\qty{192.7(4)}{\angstrom}\) with RMS surface roughness of \(\qty{6.8(4)}{\angstrom}\). Errors are reported as \(2\sigma\) from the covariance matrix.}\label{tab:4.1}
  \begin{tabular}{l S[table-format=2.1(1)] S[table-format=1.2(1)] S[table-format=2.1(1)]}
    \toprule
    Region & {\(\gamma\) (\unit{\degree})} & {\(\rho\) (\unit{g.cm^{-3}})} & {\(\tau\) (\unit{\angstrom})} \\
    \midrule
    Surface & 58.7(6) & 2.0(1) & 13(1) \\
    Bulk & 69.6(2) & 1.64(1) & {---} \\
    Interface & 38(1) & 0.87(7) & 14.9(1) \\
    \bottomrule
  \end{tabular}
\end{table}

The recovered structure is plotted in \cref{fig:4.4}b and~d with fit parameters in \cref{tab:4.1}. The bulk tilt is consistent with the film averaged orientation from AR-NEXAFS (\suppfig{5}). At the buried interface, the profile descends to a lower density, face-on orientation, while at the free surface the orientation relaxes to a lower angle and a slightly higher density than bulk. This profile is consistent with a second film grown under identical conditions with a different total thickness (\suppfig{10}). The comparison and consistency between independent films, suggests these parameters may be intrinsic to the growth conditions of the ZnPc system and dictated directly by the free energy potentials at the vacuum and silica interfaces.

\Cref{fig:4.4}c presents a molecular level, real space reconstruction of the film through data fusion of our RSoXS analysis, XRD lattice spacings (Figs.~S11 and~S12), and a force-field model of the molecule. Notably, the entire film is only \({\sim}\)15 monolayers thick with the substrate surface energy affecting a region equivalent to a single monolayer. However, the fact that the graded interface model fits better than the discrete slab model reveals that the continuous density and orientation variation over that region represents a real structural variation resolved by this analysis. In fact, the unique tilt in the last few angstroms at the interface, can only manifest with sub-monolayer coverage of ZnPc molecules that are pinned relatively flat on the substrate. This confined feature is well below the experimental diffraction limit (\(2\pi/q_{\max} > \qty{23}{\angstrom}\)). However, scattering from such features still modify the amplitude and phase of the measured fringes. A Fisher sensitivity analysis (\suppfig{13}) determines that these sub-monolayer features are uniquely captured and constrained by our ML-DFT optical model which locks these effects as a function of energy and polarization to enable recovery of the molecular state below the diffraction limit.

\section{Surface versus Bulk Dynamics in Film Growth}\label{sec:surface-vs-bulk-dynamics-in-film-growth}

These results are direct evidence against the model of substrate patterning of bulk molecular orientation and provide evidence toward an alternative model that incorporates intermolecular interactions at the growth front. One explanation for our results is that the orientation a planar molecule adopts during deposition reflects a competition between the substrate adhesion, which favors face-on alignment and intermolecular \(\pi - \pi\) cohesion, which favors edge-on packing. The gradient then traces the crossover from a low dimensional interfacial coupling to the three-dimensional growth mechanism that sets the bulk texture. The lower density and face-on alignment at the interface indicates that the packing state here is not the relaxed herringbone but possibly a primordial face-on brickstone packing structure which typically exhibits such properties.\autocite{derrien2022interfacial} We expect that evaluating these effects on interfacial states in molecular systems will enable control of emergent phenomena at these interfaces, for example, charge generation and recombination in organic solar cells and OLED technologies\autocite{tumbleston2014influence,yokoyama2011molecular} or charge and ion transport in transistors\autocite{pei2020semiconductor} or electrochemical devices.\autocite{khan2025local,flagg2023resonant,mayer2026lumen,trigg2018self,dong2010super} The same sub-monolayer sensitivity extends directly to self-assembled monolayers, where incomplete or defective coverage during assembly is now recognized as the limiting factor for both efficiency and thermal stability in state-of-the-art perovskite solar cell devices.\autocite{jiang2025toughened} More broadly the analysis can be applied to any interface, polymer, alloy, or few-layer system whose function depends on a coverage-dependent structure that a bulk-averaged measurement cannot resolve.

The surface densification measured here, along with an orientation approaching the magic angle also indicates a distinct state from the bulk, where the higher density and disorder may be caused by the physics of surface tension during deposition. In this case, these boundary layers are not the adventitious hydrocarbon contamination nor oxidation to which film-averaged analyses and other RSoXS studies customarily assign them\autocite{ferron2023composition,fostiropoulos2019temperature} but rather potentially a disordered state before the ZnPc molecules fully crystalize. Moreover, this demonstrates that in studies characterizing structure of films deposited with varying coverages\autocite{derrien2022interfacial,yamamoto2022variation,deoteyza2006structural} do not measure the state expected in a continuously deposited film, since both substrate and free surface energetics pull the state away from the final buried states.

The new RSoXS analysis provides a depth sensitivity to buried molecular state applied in a broad range of physical conditions, allowing quantitative studies of interfacial phenomena like crystal growth, kinetic deposition, and solvent interaction. The same construction extends to molecules with additional orientational and chemical freedom, to polymers, alloys, and mixed oxidation states or orbital hybridization affecting unoccupied molecular orbitals. The analysis is, furthermore, applicable to scattering geometries beyond reflectivity, enabling local chemical, electronic, and physical analyses of inaccessible buried interfaces that govern function across material systems and nanostructures. Connecting scattering features that vary with energy and polarization, our ML-DFT optical model can not only provide these states directly but also push the resolution well beyond the traditional diffraction limit.

In summary, we introduce a resonant X-ray scattering analysis that directly interprets the quantitative physicochemical molecular state in a buried nanostructure based on an ML-refined DFT optical model. The analysis substantially improves parameter precision and resolves structure finer than the diffraction limit, so that the molecular state at an interface can be characterized directly. Applying this analysis to ZnPc films, we find that hypotheses of chemical contamination at the film surface and buried interface are rejected in favor of the molecules exhibiting a novel packing and orientational state. Introducing a structural model based on the dynamics of film growth, we reveal evidence to reject hypotheses of substrate patterning of the bulk packing state in favor of a competition between two- and three-dimensional growth that determines the interfacial evolution of the molecular states. This statistical analysis of the final buried molecular state contrasts with checkpoint studies that track a growth front that is likely only in a temporary state. We expect this non-destructive view inside buried interfaces to enable deliberate control of the interfacial states that govern function in organic devices.

\section{Methods}\label{sec:methods}

\textbf{Film growth:} Small molecule ZnPc (C32H16N8Zn sourced from Sigma Aldrich) with purity \(\geq 99.5\%\) was further purified in the physical vapor deposition (PVD) chamber (base pressure \(< 1 \times 10^{- 7}\) Torr) by gradually heating, approaching ZnPc sublimation temperature over approximately an hour until no further off gassing of contaminants were detected by the quartz crystal monitor (QCM). Thin films (\({\sim}\)\qty{20}{nm} measured by QCM) of ZnPc were deposited at \(.2\ \mathring{\mathrm{A}}/s\) on high conductivity polished Si wafers (University Wafers). Wafers were cleaved \qty{1}{cm} x \qty{3}{cm} and cleaned via sonication in deionized water, acetone and isopropyl alcohols before a 30 min UV ozone bake. ZnPc was either deposited directly onto these wafers to create the herringbone crystal structure, or onto a thin Copper Iodide (CuI, \qty{3}{nm}) patterning layer producing a face on brick-stone polymorph. Substrate temperature was either set to the ambient room temperature (24 C), by liquid nitrogen cooling (-50 C), or by thermal substrate heating (100 C). For brick-stone samples, CuI and ZnPc depositions were performed sequentially under the same vacuum, with 20 min between depositions to return to the base pressure. Films for RSoXS were stored under an inert N2 atmosphere and transported to the ALS under vacuum. Auxiliary pre characterization measurements were performed on separate films from the same batch to avoid surface contamination and oxidation prior to RSoXS.

\textbf{VASE:} Variable angle spectroscopic ellipsometry (VASE) (\suppfig{1}) was performed using a J.A. Woollam alpha-SE to pre characterize layer thickness and molecular orientation. Data was collected at three incidences (\ang{65}, \ang{70}, and \ang{75}) targeting a sample from the same PVD batch and measured using vendor supplied CompleteEase software. Optical constants where modeled using a B-spline producing similar spectra to literature.\autocite{kozlik2013optical} The results and analysis of VASE data are discussed in Supplemental Material, \suppsec{1}

\textbf{GIXRD:} Grazing incidence wide angle x-ray diffraction (GIXRD) measurements (\suppfig{2}) were performed at Washington State University on a Xenocs Xeuss 3.0 with a focused Cu X-ray source (\(\lambda = \qty{1.5418}{\angstrom}\)), under vacuum at \qty{72}{mm} sample-detector distance with incident angle set to \(\ang{0.2}\). The sample-to-detector distance and beam center were calibrated with a silver-behenate standard. 1D reduction of GIWAXS images was done using NIKA software\autocite{ilavsky2012textless} and peak positions were obtained by fitting Gaussian line shapes to the reduced profiles. The lamellar reflection lies along the \(q_{z}\) and the \(\pi - \pi\) reflection is in plane corresponding with the edge on texture of the \(\alpha\) herringbone polymorph reported for vacuum deposited ZnPc and CuPc.\autocite{hoshino2003redetermination} The results and analysis of GIXRD data are discussed in Supplemental Material, \suppsec{1}.

\textbf{AFM:} Atomic force microscopy (AFM) height maps were acquired on a Bruker Dimension Icon in tapping mode with a TESPA-V2 SAMPLE probe. Scans where processed using Gwyddion\autocite{necas2012gwyddion} to extract root-mean-square (RMS) roughness values. Film morphologies and RMS roughness values were used to pre-characterize film surface roughness and morphology RSoXS measurements. The surface topology is included in \suppfig{4} and the discussion in Supplemental Material, \suppsec{1} and \textbf{1.4}.

\textbf{NEXAFS spectroscopy:} NEXAFS spectra were collected at the Soft X-ray Spectroscopy (SXR) beamline of the Australian Synchrotron, part of ANSTO in total electron yield mode on Si wafers.\autocite{heilman2026x} Samples where grown under the same conditions Spectra were collected at five sample alignments with respect to the incident electric field polarization vector with respect to the surface normal (\(\ang{20},\ \ang{40},\ \ang{55},\ \ang{70},\ \&\ \ang{90}\)). NEXAFS was collected over the energy range of 270-\qty{350}{eV} and normalized to the energy dependent intensity fluctuation and optical loss via a direct photodiode measurement, and Au mesh drain current. Mass absorption was computed via normalization in the pre and post edge regions to the bare atom mass absorption of ZnPc as extracted from the CXRO Henke atomic scattering factors.\autocite{henke1993x}

Molecular tilt angle was extracted via the DFT cluster and refinement process. As the final step of the refinement process, the molecular polar tilt angle \(\gamma\) is fit using the TDM cluster orientations prescribed by DFT providing a measure of the molecular orientation (\suppfig{5}). This fit produces a reduces \(\chi^{2}\) goodness of fit statistic of .213 indicating very good agreement between the model and experimental data. The molecular tilt angle of \(\ang{74}\) agrees with all other characterizations of the ZnPc edge-on herringbone polymorph described in this study and is taken as the coarse-grained molecular orientation within the films measured under RSoXS (full analysis in Supplemental Material, \suppsec{3}).

\textbf{Density Functional Theory:} A time independent transition potential calculation using the StoBe library was performed on a single molecule of ZnPc. A HCH approach sets the manifold of transition energies, along with each TDM and oscillator strength, while the \(\Delta KS\ \)method is used to fix the absolute energy associated with each transition. For details see our previous work.\autocite{murcia2026quantitative} The molecular geometry was extracted from a crystallographic structure file (.cif)\autocite{mai2020data} as it produced results with a higher initial agreement to experiment compared to thermodynamically relaxed structures, or those relaxed from molecular force fields. The final calculation results are provided in Supplemental Material, \suppsec{1}.

\textbf{ML-clustering and refinement:} The ML optimized optical tensor is constructed through hierarchical agglomerative clustering to extract principal TDM clusters. Exact implementation details may be found on GitHub or through our previous work.\autocite{murcia2026quantitative,victormurcia2025wsu} TDM clusters are sequentially refined to quantitative agreement with experimental AR-NEXAFS by fitting oscillator strength, excitation energy, and TDM orientation. Optimal hyperparameters, oscillator strength threshold and overlap percentage were investigated by evaluating this best fit metric on an even grid \(OS \in \lbrack 0,\ 10\rbrack\ \&\ OVP \in \lbrack 25,75\rbrack\) with details in Supplemental Material, \suppsec{3} and highlighted in \suppfig{7}. From the family of hyperparameters that perform well on this metric, the refinement performance is additionally weighted by the adherence to the original DFT spectra and TDM cluster families such that the refinement step does not significantly alter the DFT calculation. Here we chose the best clustering parameters of \(OS = 0.25\) and \(OVP = 52.5\) with the full evaluation matrix in \suppfig{7}. In general, we find good clustering and convergence for all \(OVP \in \lbrack 50,55\rbrack\), with significantly worse performance outside of these bounds. Moreover, we find that \(OS \in \lbrack.1,8\rbrack\) performed similarly, yet parameter agreement with originating DFT degrades rapid for larger \(OS\), while the refinement never converse for \(OS < .1\).

\textbf{RSoXS collection:} Polarized resonant soft x-ray reflectivity was measured at beamline 11.0.1.2 of the Advanced Light Source, Lawrence Berkeley National Laboratory, following the procedure established by Ferron et al.\autocite{ferron2023composition} Samples are transferred from inert storage to the end station where specular reflectance is captured on the CCD area detector with thermals managed by thermos-electric cooling. The sample surface is aligned to the goniometer axis of rotation by an iterative process that refines the sample rotation and height. Incident flux was attenuated using the scatter slits together with the higher order suppressor to avoid detector damage. Because the total reflectance varies by more than the detectors dynamic range, reflectivity is captured in overlapping \(\mathbf{q}\) subsets with suppressor settings and exposure time raised at each boundary once reflectivity reaches the noise floor.

Data is collected first in s polarization and then in p polarization beginning at \qty{250}{eV}, below the carbon K edge, and stepping upward through the resonance, for 21 photon energies collected in batches of five energies with the axis of rotation alignment was reevaluated between energy batches. Upon polarization change, optics were allowed to re equilibrate to the new heat load for 30 min before alignment was recalibrated. The complete dataset was acquired within a single 12 h shift. Photon energy was calibrated against the transmission NEXAFS of a polystyrene standard at the beginning and at the end of the shift with no discernable shift, bounding the energy drift over the experiment to the width of the primary polystyrene feature.

\textbf{RSoXS reduction:} Reflectance was reduced from the two-dimensional detector images using the RSoXS-reduce package.\autocite{alsrsoxspxrreduce2026als} The specular spot was located on each frame and its intensity integrated over a fixed region of interest. An equal region displaced along the same detector row provided the background, so that subtraction removes the systematic row-wise readout offset in addition to dark current and stray light. Reflectance was then formed by dividing the background-corrected specular intensity by the incident \(I_{0}\) recorded at the start of each reflectivity profile through the averaging of five direct beam measurements.

At each energy, full multi-decade reflectivity profiles are constructed by properly scaling and joining the overlapping \(\mathbf{q}\)-subsets measured at different attenuation and exposure time. An effective \(I_{0}\) is computed from the four overlapping points, and the subsets were joined to form a singular continuous reflectivity profile for each measured energy. Reported uncertainties combine per-pixel effective \(I_{0}\) with a \(\mathbf{q}\)-dependent envelope accounting for the increasing ambiguity between specular and diffuse reflectivity at higher \(\mathbf{q}\) using the following equation:

\begin{equation}
  \sigma(E,\mathbf{q}) = a R_{\mathrm{exp}}(E,\mathbf{q}) + b\,\mathbf{q}.
  \label{eq:uncertainty-envelope}
\end{equation}

With \(a = .08\) and \(b = .5 \times 10^{- 6}\), which is a feasible approximation for experimental error under these considerations (discussion in Supplemental Material, \suppsec{5}).

\textbf{RSoXS structural and optical modeling:} Reflectivity is computed from the 4x4 transfer matrix using the \pkg{refloxide} package.\autocite{alsrsoxsrefloxide2026als} The film was modeled as three ZnPc slabs, denoted surface, bulk and interface, on a native SiO\textsubscript{2} layer and a semi-infinite Si substrate. Thickness and roughness of each ZnPc slab together with the oxide thickness and roughness give eight structural parameters. Reflection and transmission coefficients were computed at each photon energy and scattering vector with the 4x4 transfer matrix formalism matching boundary conditions at the five interfaces, with interfacial roughness entering in the Nevot-Croce sense as an exponential attenuation representing off-specular diffuse scattering from a faceted interface.\autocite{nevot1980caracterisation} Vacuum, silicon dioxide, and silicon were assigned to tabulated bare atom optical constants,\autocite{henke1993x} with an energy-independent density allowed to vary within (2, 2.8) g cm\textsuperscript{-3} to accommodate stochiometric deviation from the tabulated values. The forward model was evaluated using \pkg{refloxide}, and the objective and parameter machinery where provided by \pkg{refnx}.\autocite{nelson2019textless}

Two optical parameterizations were compared on this identical structure. The free tensor parameterization assigns each ZnPc slab an independent birefringence, dichroism, isotropic \(\beta\) and isotropic \(\delta\) at every measured energy, constraining to ranges dictated by the starting optical model, amounting to 252 optical tensor components. The DFT parameterization instead describes each slab energy-independently by a molecular orientation and a density, with one energy offset term common to all layers absorbing the calibrated difference between the NEXAFS and RSoXS instruments, amounting to seven parameters for all three slabs. A detailed parameter accounting and comparison with prior art is discussed in Supplemental Material, \suppsec{5}.

Experimental corrections were applied per energy and per polarization as a normalization scale factor and an angle-of-incidence offset, following Ferron et al.,\autocite{ferron2023composition} adding four terms per energy at 84 in total. The refinement therefore varies 101 parameters for the DFT model and 350 for the free tensor model. Bound on the correction terms, and the correlations they induce are described in Supplemental Material, \suppsec{5}.

All energies and polarizations are co-refined against a single global objective, so that the structural parameters are shared across the dataset and only the optical and correction terms carry energy indices. The log-likelihood assumed Gaussian errors at the uncertainties propagated though reduction, and all varying parameters were assigned uniform priors over their stated bounds. The maximum likelihood parameter set was located by differential evolution over the full bounded space, in parallel with deferred population updating and a local polish step using the L-BFGS-B algorithm, avoiding dependency on initial conditions. Reported parameter uncertainties are then taken from the inversion of the Hessian evaluated in the final polish step. Model comparison uses reduced chi square together with the Akaike and Bayesian information criteria.

\textbf{Unit cell construction:} Unit cells used for molecular dynamics simulations are constructed though fusion of GIXRD \(d\)-spacings with literature unit cells. The rigid ZnPc molecule originates from the same \(\beta\)-ZnPc (\suppfig{12}a) crystallographic results used for DFT simulations.\autocite{mai2020data} Lennard-Jones parameters are GAFF2, with UFF for the Zn coordination center. Charges are fitted using the B3LYP/def2-SVP electrostatic potentials and used to set the electrostatic repulsion between molecules.

The \(\beta\)-ZnPc stable polymorph with a known structure serves only as the validation target. A suitable unit cell was constructed from the chemically similar \(\alpha\)-CuPc by placing rigid ZnPc molecules within the CuPc unit cell. Lattice parameters \textbf{a} and \textbf{c} are set by the lamellar spacing (002), with \textbf{b} set by the \(\pi - \pi\) packing. The molecular tilt angle is scanned and compared to the GIXRD patterns though a fiber-texture forward model, with other unit cell angles set by minimizing the lattice energy (UFF+ESP). The resulting lattice parameters are tabulated in \supptab{3} with the resulting crystallographic (.cif) structure available and visualized in \suppfig{12}b. The final unit cell is then thermally relaxed by rigid-body NPT at \qty{300}{K}, with independent barostats on each axis such that the force fields set the lattice and herringbone arrangements. Over this range the minimum enthalpy cell produces a tilt of \(\ang{72(1)}\) and a density of \(1.657\) g/cm\textsuperscript{3} against \(\ang{69.6(2)}\) and \(1.64(1)\) g/cm\textsuperscript{3} from RSoXS and reproduces the lamellar and (002) spacings within 1.5\%. The final relaxed unit cell is visualized in \suppfig{12}c used in the remaining MD simulations.

\textbf{Real space molecular visualization:} Molecular packing in \cref{fig:4.4}c is a molecular dynamics (MD) reconstruction of the film. It is built from the fitted depth profiles of \cref{fig:4.4}a and~b and relaxed in LAMMPS.\autocite{thompson2022lammps} ZnPc is treated as a rigid body in the molecular frame used for the DFT calculations. Lennard-Jones parameters are taken from GAFF2, with UFF for Zn, and atomic charges are fitted to the B3LYP/def2-SVP electrostatic potential. This force field reproduces the \(\beta\)-ZnPc unit cell at \qty{100}{K} to within 1.5\% on each axis (SI).

The film is constructed for pure visualization by data fusion with the fitted depth profiles, GIXRD, and the fitted electrostatic potentials. The molecular frame is the same used for the DFT calculations (C32H16N8Zn, \(M = 577.9\) g/mol). It's van der Waals envelope gives an in-plane diameter of 14.5 \(\mathring{\mathrm{A}}\) and a face-to-face thickness of \(\qty{3.6}{\angstrom}\), and the \(\pi\)-stacking separation is taken as \(\qty{3.4}{\angstrom}\). At each depth \(z\ \)the molecular tilt is \(\gamma(z)\) from the orientation profile, and the area per molecule in the drawing plane is \(M/(\rho(z)N_{A}t)\), with \(\rho(z)\) the fitted density profile and \(t = \qty{14.4}{\angstrom}\) the extrusion thickness of the slab that the projection represents. The packing density of the molecules in the illustration is therefore self-consistent with the density and orientation gradients from \cref{fig:4.4}b and~d. The projected cell is orthogonal, carrying a lamellar period along \(z\) and a column period in plane whose ratio is fixed by the molecular diameter, the \(\pi\) stacking separation and \(\sin^{2}\left( \gamma(z) \right)\). The lamellae period is modified by the interdigitation between neighboring lamellae, set by a dimensionless factor \(f = 0.96\), so adjacent layers can overlap by \(4\%\) of the molecular diameter. This together with the herringbone splay and positional and orientational disorder are set only for visual continuity and are neither empirical nor derived factors.

The depth dependent lamellae and \(\pi\) d-spacings are compared with diffraction orders derived from experimental GIWAXS measurements in \suppfig{11}. Because both the cell area and its aspect ratio vary with depth, each z monolayer contributed its own pair of spacings rather than a single value, so the predicted diffraction positions are a taken as the maxima of these depth distributions, weighted by the local density and the Debye-Waller factor built from the drawn positional disorder. The lamellar and \(\pi\) spacing periods agree with the measured peak centers to around \(2\%\), treated as a secondary consistency check of the molecular visualization.

\textbf{Data Availability:} NEXAFS data that support the findings of this study have been deposited in X-ray Atlas (\url{https://xrayatlas.wsu.edu}) with the accession code \href{http://xrayatlas.wsu.edu/d/193bc0c5}{193bc0c5} and are archived at Zenodo (\url{https://doi.org/10.5281/zenodo.22742990}).\autocite{heilman2026x} RSoXS reflectivity curves, GIWAXS, non-resonant XRR, VASE, and AFM topography together with the fitted model parameters that support the findings of this study are available from the corresponding author upon reasonable request.

\textbf{Code availability:} The source code used for reflectivity analysis may be found on GitHub (\url{https://github.com/wsu-carbon-lab/refl-analysis}).\autocite{heilman2026refl} The code for the RSoXS 1D to 2D data reduction can be found on GitHub (\url{https://github.com/als-rsoxs/pxr-reduce}).\autocite{alsrsoxspxrreduce2026als} The source code used to compute reflectivity and fit RSoXS data may be found on GitHub (\url{https://github.com/als-rsoxs/refloxide}).\autocite{alsrsoxsrefloxide2026als} The source code used for the DFT clustering and optical model generation may be found on GitHub (\url{https://github.com/wsu-carbon-lab/dft-learn}).\autocite{victormurcia2025wsu}

\textbf{Acknowledgements:} This research used resources of the Advanced Light Source, which is a DOE Office of Science User Facility under contract no. DE-AC02-05CH11231. Part of this work was carried out on the Soft X-ray Spectroscopy (SXR) beamline of the Australian Synchrotron, part of ANSTO. The authors also acknowledge the National Energy Research Scientific Computing Center (NERSC) provide supercomputer resources to this work for the DFT calculations, and RSoXS fits. The authors also acknowledge Dr. Thomas Ferron, Dr. Cheng Wang, and Dr. Kas Andrle form the Advanced Light Source of Lawrence Berkeley National Laboratory for their training and guidance in data collection, reduction, and fitting analysis. The authors also thank Dr. Arda Gozen from the School of Mechanical and Materials Engineering of Washington State University for AFM support on measuring soft nanoscale films. Fable 5 from Anthropic was used for minor copy editing and code generation under the ALS computing group's team plan. The authors also thank Dr. Daniel Dougherty from North Carolina State University and Dr. Dean DeLongchamp from the National Institute of Standards and Technology for their reading and consultation of the manuscript.
